\chapter{The Simple Sentence}\label{ch:simple}

% ============================================================
% STATUS: SCAFFOLD — largely new content; some material available in
% archive/2020/chapters/07-clause-structure.tex.
% Key theoretical framing: topic-comment structure is the primary
% organizational principle of the Iridian clause. The split-ergative
% alignment can be recast as: the absolutive is typically the topic;
% the ergative demotes the agent away from topic status.
% ============================================================

\section{Topic-comment structure and information structure}\label{sec:topic-comment}

% ============================================================
% The basic Iridian clause is organized around a topic-comment structure.
% The topic (what the sentence is "about") is unmarked (absolutive) and
% typically appears sentence-initially. The comment (what is said about
% the topic) follows. The alignment system is closely tied to this:
%   - Perfective/retrospective: P is the topic (ABS), A is demoted (ERG)
%   - Nominative aspects: A/S is the topic (ABS), P is non-topic (OBL)
% This means aspect choice is partly pragmatic — the speaker chooses the
% aspect that places the intended topic in the unmarked (absolutive) position.
% Antipassive demotes P and restores A to topic/ABS position.
% ============================================================

The basic Iridian clause organizes information around a topic-comment structure. The \textbf{topic} --- the entity the sentence is about --- is expressed in the absolutive case and typically occupies the initial position in the clause. The \textbf{comment} --- what is predicated of the topic --- follows. This organization is reflected in and reinforced by the split-ergative alignment: in the perfective and retrospective aspects, the patient is promoted to the absolutive (becoming the topic), while in the nominative aspects, the agent occupies the absolutive (topic) position.

Aspect choice is therefore partly a pragmatic decision: the speaker selects the aspect that places the discourse-salient argument in the absolutive position.

% [Examples to be added]

\section{Word order}\label{sec:word-order}

The canonical word order in Iridian is SOV (subject-object-verb), where the finite verb is clause-final. This order is relatively rigid for the finite verb itself, which must appear last. The ordering of nominal arguments before the verb is more flexible and is governed by information structure rather than strict grammatical requirements. Topic-initial and focus-final tendencies apply within the pre-verbal domain.

% [Examples: SOV default, pragmatic reordering, focus movement]

\section{The copula and existential constructions}\label{sec:copula}

Iridian distributes predication across three structurally and semantically distinct constructions. Where many European languages rely on a single copular verb for all such functions, Iridian treats equative identity, marked stative or inchoative predication, and existential predication separately. All overt predicative verbs are clause-final and inflect for the full aspect system; each paradigm is strongly suppletive and morphologically irregular.

\subsection{Equative and class-membership clauses: the zero copula}\label{sec:zero-copula}

The default pattern for equative clauses --- those asserting identity, role, or class membership --- is zero copula in the indicative. The predicative complement appears in the direct case in clause-final position; no overt verb is present. This is the stylistically neutral, unmarked construction.

\pex
\a
\begingl
\gla Marek doktor.//
\glb Marek.\Dir{} doctor.\Dir{}//
\glft \trsl{Marek is a doctor.}//
\endgl
\a
\begingl
\gla Iridian malno.//
\glb Iridian.\Dir{} language.\Dir{}//
\glft \trsl{Iridian is a language.}//
\endgl
\xe

Inclusive and class-membership predications --- of the type \textit{dogs are animals} --- follow the same zero-copula pattern and carry no inference beyond simple predication. The predicative complement always appears in the direct case regardless of which case would otherwise be assigned in the same position.

\subsection{The marked copular verb \ird{neh-}: becoming and inchoative predication}\label{sec:neh}

The root \ird{neh-} provides the overt copular verb. Its paradigm is strongly suppletive and morphologically irregular; principal forms are given in Table~\ref{tab:neh-paradigm}.

\begin{table}[ht]
	\footnotesize\sffamily
	\caption{Principal forms of \ird{neh-} (marked copula).}\label{tab:neh-paradigm}
	\medskip
	\begin{tabularx}{0.72\textwidth}{L l}
		\toprule
		{\sc form} & {\sc iridian} \\
		\midrule
		Continuous          & \ird{neha} \\
		Progressive         & \ird{nehime} \\
		Contemplative       & \ird{nehach} \\
		Perfective          & \ird{nek} \\
		Retrospective       & \ird{neaní} \\
		Sequential converb  & \ird{nesu} (colloquial \ird{nu}) \\
		Simultaneous converb & \ird{nešěc} \\
		\bottomrule
	\end{tabularx}
\end{table}

Since neutral equative predication is expressed by zero copula (§\,\ref{sec:zero-copula}), overt \ird{neh-} is semantically marked. It is not used for simple statements of identity or class membership. Instead it signals \textbf{becoming}, \textbf{entering a state}, or \textbf{transition}: the subject is understood as moving from one classification or condition into another.

\pex
\a
\begingl
\gla Marek doktor neha.//
\glb Marek.\Dir{} doctor.\Dir{} \Cop{}-\Cont{}//
\glft \trsl{Marek is becoming a doctor.}//
\endgl
\a
\begingl
\gla Ša byl kapitan nek.//
\glb \Dem{}.\Prox{} child.\Dir{} captain.\Dir{} \Cop{}-\Pf{}//
\glft \trsl{This child became captain.}//
\endgl
\xe

In the continuous aspect, \ird{neh-} may describe a newly entered or recently changed state, with the inchoative flavour modulated by context; a purely stable state assertable without implication of transition calls for zero copula instead. The predicative complement of \ird{neh-} stands in the direct case, as with zero-copula equatives.

The sequential converb \ird{nesu} --- colloquially \ird{nu} --- is a suppletive form established independently of the regular aspect paradigm. Its grammaticalization as a discourse particle is described in Section~\ref{sec:seq-copula} of Chapter~\ref{ch:complex}.

\subsection{The existential verb \ird{jec-}: existence, presence, and having}\label{sec:jec}

The existential verb \ird{jec-} is a true clause-final verb asserting the existence or presence of a referent. It inflects for all aspects; its paradigm is strongly suppletive and morphologically irregular, with the perfective and retrospective built on a distinct suppletive root. Principal forms are given in Table~\ref{tab:jec-paradigm}.

\begin{table}[ht]
	\footnotesize\sffamily
	\caption{Principal forms of \ird{jec-} (existential verb).}\label{tab:jec-paradigm}
	\medskip
	\begin{tabularx}{0.72\textwidth}{L l}
		\toprule
		{\sc form} & {\sc iridian} \\
		\midrule
		Continuous          & \ird{jeca} \\
		Progressive         & \ird{jecime} \\
		Contemplative       & \ird{jecach} \\
		Perfective          & \ird{žek} \\
		Retrospective       & \ird{žaní} \\
		Sequential converb  & \ird{žu} \\
		\bottomrule
	\end{tabularx}
\end{table}

The division of labour between \ird{jec-} and the zero-copula construction is loosely parallel to the Japanese distinction between the existential verbs \textit{iru} and \textit{aru} on the one hand and the equative \textit{da}/\textit{desu} on the other: \ird{jec-} covers existence and presence; zero copula covers identity and class membership.

\paragraph{Existence and presence.} The primary use of \ird{jec-} is to assert that a referent exists or is present. The entity in question stands in the direct case.

\pex
\a
\begingl
\gla Byl jeca.//
\glb child.\Dir{} \Exst{}-\Cont{}//
\glft \trsl{There is a child. / A child is present.}//
\endgl
\a
\begingl
\gla Kroumaště na po zma jeca pivo.//
\glb refrigerator.\Ind{} in still few \Exst{}-\Cont{} beer.\Dir{}//
\glft \trsl{There is still some beer left in the refrigerator.}//
\endgl
\xe

\paragraph{Location.} \ird{Jec-} is also used when locating an entity: the entity appears in the direct case and the location is expressed by an indirect noun phrase followed by a postposed clitic adposition.

\pex
\a
\begingl
\gla Marek dumě na jeca.//
\glb Marek.\Dir{} house.\Ind{} in \Exst{}-\Cont{}//
\glft \trsl{Marek is at home.}//
\endgl
\xe

\paragraph{Having.} Possession is expressed through \ird{jec-} with an oblique possessor: `I have X' is construed as `there is X at me', with the possessor in the indirect case and the possessed item in the direct.

\pex
\a
\begingl
\gla Bylam jeca tóm.//
\glb child.\Ind{} \Exst{}-\Cont{} book.\Dir{}//
\glft \trsl{The child has a book. (Lit.\ At the child there is a book.)}//
\endgl
\xe

\subsection{The negative existential verb \ird{zad-}}\label{sec:zad}

The root \ird{zad-} is the lexically negative counterpart of \ird{jec-}: it asserts the non-existence or absence of a referent. Because \ird{zad-} is inherently negative, it does not combine with the standard negation prefix \ird{zá-}. Its paradigm is given in Table~\ref{tab:zad-paradigm}.

\begin{table}[ht]
	\footnotesize\sffamily
	\caption{Principal forms of \ird{zad-} (negative existential verb).}\label{tab:zad-paradigm}
	\medskip
	\begin{tabularx}{0.72\textwidth}{L l}
		\toprule
		{\sc form} & {\sc iridian} \\
		\midrule
		Continuous          & \ird{zada} \\
		Progressive         & \ird{zadime} \\
		Contemplative       & \ird{zadach} \\
		Perfective          & \ird{žatek} \\
		Retrospective       & \ird{žataní} \\
		Sequential converb  & \ird{zadu} \\
		\bottomrule
	\end{tabularx}
\end{table}

\pex
\a
\begingl
\gla Pivo zada.//
\glb beer.\Dir{} \NegExst{}-\Cont{}//
\glft \trsl{There is no beer.}//
\endgl
\a
\begingl
\gla Marek zada.//
\glb Marek.\Dir{} \NegExst{}-\Cont{}//
\glft \trsl{Marek is not here. / Marek is absent.}//
\endgl
\xe

The plural particle \ird{ne} cannot be combined with a \ird{zad-} clause; see Section~\ref{sec:definiteness} of Chapter~\ref{ch:nouns}.

\section{Questions}\label{sec:questions}

\subsection{Yes-no questions}

% [To be written: polar question formation — intonation, question particle]

\subsection{Content questions}

Content questions (\textit{wh}-questions) require the interrogative to occupy the topic (initial) position of the clause. \textit{Wh}-fronting is obligatory in Iridian; in-situ interrogatives are not grammatical in standard Iridian.

% [Port from 2020 archive §\,sec:content-questions; update case labels]

\subsection{Tag, echo, and indirect questions}

% [To be written]

\section{Negation in clausal context}\label{sec:clausal-negation}

At the clausal level, negation is expressed by the verbal prefix \ird{zá-} (reduced to \ird{za-} in unstressed environments). The prohibitive mood suffix \ird{-éma} handles negation of commands. The negative existential verb \ird{zad-} handles negation of existential constructions; because it is lexically negative, it does not take the prefix \ird{zá-} (see §\,\ref{sec:zad}). Interaction with the perfective aspect (negated perfective = `should have done but didn't') is discussed in Chapter~\ref{ch:verbs}.

% [Full treatment to be added]

\section{Comparative and superlative constructions}\label{sec:comparison}

% ============================================================
% The 2020 archive described a comparative construction using the agentive
% case (now: transitive case) for the standard of comparison, and a comparative
% particle. This needs to be updated to the new case terminology.
% ============================================================

% [Port and revise from 2020 archive; update "agentive of comparison" to
%  "transitive case in comparison"]
